\section{本章小结}
\label{sec:pspace-summary}

本章把第一部分分散的模型放进同一个框架。一个概率模型由三部分组成：样本空间$\Omega$列出可能的结果，事件族$\mathcal{F}$规定可以问概率的结果集合，概率测度$\mathbb{P}$给每个事件指定概率。事件族之所以要单独指定，是因为连续位置上无法给全部子集一致地赋予长度，也因为同一组结果在不同的观察者那里能回答的问题不同。它对补集和可数并封闭，因而是$\sigma$代数；在有限样本空间上，它就是一个划分及其各种并。

概率是总质量为一的测度。可数可加性推出概率沿单调事件列连续，由此得到分布函数的右连续性、两端极限和跳跃等于单点概率；分布函数又唯一确定实轴上的概率测度，任何满足三条性质的函数都对应一个分布，Cantor分布也在其中。单点概率为零与总概率为一并不矛盾，因为可加性只对可数个事件成立。零概率事件可以发生，几乎必然的陈述可以合并可数次，但不能合并不可数次；完备化把零概率事件的子集补为事件，使零概率修改不破坏可测性，却不改变任何原有概率，也不能使零概率条件变得可以相除。

随机对象是可测映射，可测性保证由观测引出的“取值落在某集合中”都是事件。实值、向量、矩阵和有限随机图都是它的特例，可测函数的复合仍是随机对象。观测生成的$\sigma$代数记录了它能区分的事件，对观测再作变换只会减少信息。随机对象的分布是概率在它下面的推前，变换后的分布由原像规则直接算出；分位数变换说明任何一维分布都可以由一个均匀随机点实现。处处相等、几乎必然相等与同分布依次减弱，边缘分布不能确定两个变量怎样配对，差与和的分布需要联合分布。

独立是信息族之间的乘法公式。它只需在对交封闭的生成族上检验，所以联合分布函数分解为乘积就足以推出独立；独立变量分组后的函数仍然独立。乘积测度造出了具有给定分布的独立对，独立等价于联合分布是边缘分布的乘积；木棍上一个均匀点的二进制数字是一列独立的公平比特，把这些数字拆组再作分位数变换，就得到具有任意分布的独立序列。

期望是随机变量对概率测度的积分：非负变量先用按取值分层的简单变量逼近，有符号变量再拆成正部与负部。线性、保序、单调收敛与支配收敛、分布积分公式、尾积分公式和独立乘积的期望，都是卷四积分理论的概率形式；第一部分的求和、密度积分和混合计算是它们的特例，等待时间的期望$0.7$与Cantor变量的方差$1/8$由此算出。

本章没有处理的是给定信息之后的平均。$\sigma(X)$描述了观测$X$提供的信息，下一章研究固定这部分信息、在每组已知结果内求平均得到的随机变量，即条件期望，并说明在观测值本身概率为零时怎样保留“组内平均”的含义。
